\documentclass[11pt]{article}
\usepackage{amsmath, amssymb}
\usepackage{geometry}
\geometry{margin=1in}
\usepackage{enumitem}



\title{Midterm Review}
\author{Xiao Zhang}
\date{February 2026}
\begin{document}
\maketitle


\section*{Part I. Basic Statistical Concepts and Data Collection}

\begin{itemize}[leftmargin=*]
\item Define and distinguish between \textbf{statistics}, \textbf{probability}, \textbf{population}, \textbf{sample}, \textbf{parameter}, and \textbf{statistic}.
\item Distinguish among common \textbf{sampling methods}:
\begin{itemize}
\item Simple random sampling
\item Stratified sampling
\item Cluster sampling
\item Systematic sampling
\item Convenience sampling
\end{itemize}
\item Identify different \textbf{types of variables}:
\begin{itemize}
\item Quantitative (discrete vs. continuous)
\item Qualitative / Categorical (binary, nominal, ordinal)
\end{itemize}
\item Understand principles of \textbf{ethical data collection and experimentation}, including informed consent, random assignment, and avoiding harm.
\item Recognize common \textbf{sources of bias}: selection bias, response bias, nonresponse bias, measurement bias, and undercoverage.
\end{itemize}

\section*{Part II. Descriptive Statistics and Data Visualization}

\begin{itemize}[leftmargin=*]
\item Interpret common \textbf{data displays}: histograms, boxplots, scatterplots, bar charts, and pie charts.
\item Calculate and interpret \textbf{measures of center}: mean, median, and mode.
\item Use sigma notation and correctly apply the sample mean formula:
\[
\bar{x} = \frac{1}{n} \sum_{i=1}^n x_i.
\]

\item Calculate and interpret \textbf{measures of spread}: range, interquartile range (IQR), variance, and standard deviation.

\item Use and interpret the formulas:
\[
s^2 = \frac{1}{n-1} \sum_{i=1}^n (x_i - \bar{x})^2,
\qquad
s = \sqrt{s^2}.
\]

\item Understand distribution \textbf{skewness} and the relationship between mean and median.

\item Understand the effect of \textbf{linear transformations} on measures of center and spread.

\section*{Part III. Probability Concepts and Counting Methods}

\subsection*{Basic Probability Rules}

\begin{itemize}[leftmargin=*]
\item Understand basic probability terminology: random experiment, sample space, event, probability, and complement.
\item Distinguish between \textbf{independent} and \textbf{mutually exclusive} events.
\item Apply the addition rule:
\[
P(A \cup B) = P(A) + P(B) - P(A \cap B).
\]

\item Apply the multiplication rule:
\[
P(A \cap B) = P(A)\,P(B \mid A).
\]

\item Calculate conditional probabilities:
\[
P(A \mid B) = \frac{P(A \cap B)}{P(B)}.
\]

\item Apply the Law of Total Probability:

If $\{B_1,\dots,B_k\}$ is a partition of the sample space with $P(B_i)>0$, then
\[
P(A) = \sum_{i=1}^k P(A \mid B_i)\,P(B_i).
\]

\item Use contingency tables to compute marginal, joint, and conditional probabilities.

\item Determine whether two events are independent using probability calculations.
\end{itemize}


\subsection*{Counting Techniques}

\begin{itemize}[leftmargin=*]
\item Recognize when counting methods are needed to compute probabilities.
\item Distinguish between \textbf{permutations} (order matters) and \textbf{combinations} (order does not matter).

\item Apply permutation formulas:
\[
{}_nP_r = \frac{n!}{(n-r)!}.
\]

\item Apply combination formulas:
\[
{}_nC_r = {n \choose r} = \frac{n!}{r!(n-r)!}.
\]

\item Identify real-world scenarios where permutations or combinations are appropriate.
\end{itemize}

\section*{Part IV. Discrete Random Variables}

\begin{itemize}[leftmargin=*]
\item Understand the definition of a \textbf{random variable} and how to define one in context.

\item Distinguish between discrete and continuous random variables.

\item Recognize and apply the \textbf{binomial distribution}:
\[
P(X = k) = {n \choose k} p^k (1-p)^{n-k}, 
\qquad 
E(X)=np, 
\qquad 
SD(X)=\sqrt{np(1-p)}.
\]

\item Recognize and apply the \textbf{Poisson distribution}:
\[
P(X = k) = \frac{\lambda^k e^{-\lambda}}{k!}, 
\qquad 
E(X)=\lambda, 
\qquad 
Var(X)=\lambda.
\]

\item Choose an appropriate discrete distribution for a given scenario.

\item Apply linear transformations of random variables:
\[
E(aX+b)=aE(X)+b, 
\qquad 
Var(aX+b)=a^2Var(X).
\]
\end{itemize}

\end{document}
